% ============================================================
% Вспомогательные граничные данные и невозможность внутренней когерентности
% в подходе конечного нерва к спектральной каузальной теории
% ============================================================
\documentclass[11pt,a4paper]{article}

\usepackage{fontspec}
\defaultfontfeatures{Ligatures=TeX}
\usepackage{polyglossia}
\setdefaultlanguage{russian}
\setotherlanguage{english}
\setmainfont{Times New Roman}
\setsansfont{Arial}
\setmonofont{Courier New}
\newfontfamily\cyrillicfont{Times New Roman}
\newfontfamily\cyrillicfontsf{Arial}
\newfontfamily\cyrillicfonttt{Courier New}
\usepackage{amsmath,amssymb,amsthm}
\usepackage{mathtools}
\usepackage{enumitem}
\usepackage{booktabs}
\usepackage{tikz}
\usetikzlibrary{calc,arrows.meta,positioning,decorations.pathreplacing,backgrounds}
\usepackage[margin=2.5cm]{geometry}
\usepackage{hyperref}

\hypersetup{
  colorlinks=true,
  linkcolor=blue!60!black,
  citecolor=blue!60!black,
  urlcolor=blue!60!black,
  pdfauthor={Давид Алфёров},
  pdftitle={Вспомогательные граничные данные и невозможность внутренней когерентности в подходе конечного нерва к спектральной каузальной теории},
  pdfsubject={Спектральная каузальная теория, граничные операторы конечного нерва}
}

%%% Окружения теорем
\theoremstyle{plain}
\newtheorem{theorem}{Теорема}[section]
\newtheorem{proposition}[theorem]{Утверждение}
\newtheorem{lemma}[theorem]{Лемма}
\newtheorem{corollary}[theorem]{Следствие}

\theoremstyle{definition}
\newtheorem{definition}[theorem]{Определение}
\newtheorem{example}[theorem]{Пример}
\newtheorem{construction}[theorem]{Конструкция}

\theoremstyle{remark}
\newtheorem{remark}[theorem]{Замечание}

%%% Custom commands
\newcommand{\NN}{\mathbb{N}}
\newcommand{\ZZ}{\mathbb{Z}}
\newcommand{\cU}{\mathcal{U}}
\newcommand{\cN}{\mathcal{N}}
\newcommand{\cT}{\mathcal{T}}
\newcommand{\swapAB}{\mathrm{swap}_{01}}
\newcommand{\swapAC}{\mathrm{swap}_{02}}
\newcommand{\swapBC}{\mathrm{swap}_{12}}
\newcommand{\chosenface}{\varphi}

\begin{document}

\title{Вспомогательные граничные данные и невозможность внутренней когерентности\\
  в подходе конечного нерва к спектральной каузальной теории}

\author{Давид Алфёров\\[4pt]
  \small\textit{davidich.alfyorov@gmail.com}}

\date{\today}

\maketitle

% ============================================================
\begin{abstract}
Мы строим граничные операторы на нерве конечного покрытия с использованием
вспомогательных данных ориентации и доказываем, что совместимые данные
порождают цепной комплекс $C_2 \xrightarrow{\partial_2} C_1 \xrightarrow{\partial_1} C_0$
со свойством $\partial_1 \circ \partial_2 = 0$ и корректно определённой
первой гомологией~$H_1$.
Мы устанавливаем, что совместимые данные ориентации всегда существуют,
хотя их построение неканонично. Мы доказываем, что условие когерентности,
обеспечивающее совместимость, допускает несколько эквивалентных формулировок ---
от попарных законов ветвящих рёбер и выбранных граней до чисто треугольных
условий перекрытия.
Затем мы доказываем основной результат: для любой пары 2-симплексов,
имеющих общее граничное ребро, два набора глобальных свидетельских данных,
совпадающих всюду, кроме единственной транспозиции вершин на одном симплексе,
не могут одновременно удовлетворять глобальной когерентности; более того,
если один из них удовлетворяет когерентности, то другой неизбежно содержит
локальный конфликт на общем ветвящем ребре. Это показывает, что условие
когерентности не определяется симплициальным носителем нерва и что любой
подход к внутреннему извлечению гомологических данных должен использовать
структурную информацию, выходящую за рамки неупорядоченных симплициальных
данных нерва. Все результаты формализованы и машинно-верифицированы
в Lean~4 без использования \texttt{sorry} или \texttt{admit}.
\end{abstract}

\medskip
\noindent\textbf{Ключевые слова:} конечный нерв, граничные операторы,
симплициальная гомология, препятствие когерентности, спектральное действие,
формальная верификация

\tableofcontents

% ============================================================
\section{Введение}
\label{sec:intro}
% ============================================================

Принцип спектрального действия~\cite{Chamseddine:1996zu,Chamseddine:2006ep}
обеспечивает геометрическое происхождение лагранжиана Стандартной модели,
связанного с гравитацией, путём кодирования классического действия в
спектре оператора Дирака на некоммутативной геометрии. Центральная открытая
проблема в спектральном подходе к квантовой гравитации состоит в том,
несут ли дискретные структуры, возникающие в конечных приближениях
пространства-времени --- каузальные множества~\cite{Bombelli:1987aa,
Sorkin:2003bx}, конечные спектральные тройки~\cite{Connes:1994yd,
vanSuijlekom:2014pya} или конечные покрытия континуума --- достаточную
внутреннюю структуру для восстановления граничных и гомологических данных,
необходимых для корректно определённого спектрального действия.

В рамках спектральной каузальной теории
(СКТ)~\cite{Alfyorov:2026ff,Alfyorov:2026ss,Alfyorov:2026ct}
этот вопрос принимает следующую конкретную форму.
Пусть задано конечное покрытие~$\cU$ пространства~$X$ конечным
индексным множеством~$I$; тогда нерв~$\cN(\cU)$ наследует естественную
структуру абстрактного симплициального комплекса на уровне $0$-, $1$-
и $2$-симплексов. Определение граничных операторов на этом комплексе,
удовлетворяющих свойству цепного комплекса
$\partial_1 \circ \partial_2 = 0$, требует \emph{данных ориентации}:
выбора упорядочения вершин для каждого симплекса. Вопрос заключается в том,
могут ли такие данные быть получены \emph{внутренне} --- то есть
из одной лишь комбинаторной структуры нерва --- или же они неизбежно
требуют внешней информации.

\medskip

\noindent\textbf{Результаты.} Мы доказываем два дополняющих друг друга
результата.

\smallskip
\noindent\emph{(I) Положительный результат} (теоремы~\ref{thm:chain-complex}
и~\ref{thm:existence}). Если заданы вспомогательные данные ориентации,
удовлетворяющие \emph{условию совместимости}, то индуцированные
граничные операторы образуют цепной комплекс с~$\partial_1 \circ \partial_2 = 0$,
и первая группа гомологий~$H_1 = Z_1 / B_1$ корректно определена.
Такие совместимые данные всегда существуют, хотя их построение требует
неканонического выбора упорядоченных свидетелей на каждом $2$-симплексе.

\smallskip
\noindent\emph{(II) Отрицательный результат} (теорема~\ref{thm:main-obstruction}).
Условие когерентности, обеспечивающее совместимость, не определяется
симплициальным носителем нерва. А именно, для любой пары $2$-симплексов,
имеющих общий граничный $1$-симплекс, два набора глобальных свидетельских
данных, совпадающих всюду, кроме единственной транспозиции вершин на
одном $2$-симплексе, не могут одновременно удовлетворять глобальной
когерентности. Если один из них когерентен, то другой содержит
локальный конфликт на общем ветвящем ребре.

\smallskip
В совокупности эти результаты показывают, что граничные операторы на
конечных нервах могут быть построены из вспомогательных данных, однако
когерентность, необходимая для их объединения в цепной комплекс,
\emph{зависит от выборов, выходящих за рамки неупорядоченной
симплициальной структуры}. В контексте спектральной геометрии это означает,
что любое внутреннее восстановление гомологических данных из дискретного
слоя должно привлекать более богатую структурную информацию --- такую как
частичный порядок, метрику или спиновую структуру, --- которая не
содержится в одном лишь симплициальном носителе нерва.

\medskip

\noindent\textbf{Формальная верификация.} Все определения и теоремы
в данной работе формализованы в Lean~4 с использованием библиотеки
Mathlib4~\cite{Mathlib2024}. Полная формализация компилируется
без~\texttt{sorry}, \texttt{admit} или объявлений~\texttt{axiom}
и включает 63 модуля с более чем 260 утверждениями уровня теорем.
Подробности приведены в приложении~\ref{app:lean}.

\medskip

\noindent\textbf{Структура работы.}
Раздел~\ref{sec:nerve} вводит конечные покрытия и их нервы.
Раздел~\ref{sec:orientation} определяет данные ориентации и граничные
операторы.
Раздел~\ref{sec:chain-complex} доказывает свойство цепного комплекса
при наличии совместимости.
Раздел~\ref{sec:existence} устанавливает неканоническое существование
совместимых данных.
Раздел~\ref{sec:coherence} развивает эквивалентные формулировки условия
когерентности.
Раздел~\ref{sec:obstruction} содержит основной отрицательный результат.
Раздел~\ref{sec:discussion} обсуждает следствия и открытые вопросы.

% ============================================================
\section{Конечные покрытия и их нервы}
\label{sec:nerve}
% ============================================================

\begin{definition}[Конечное покрытие]
\label{def:cover}
Пусть $X$ --- множество, а $I$ --- конечное множество. \emph{Конечное
покрытие} множества~$X$ множеством~$I$ --- это семейство
$\cU = (U_i)_{i \in I}$ подмножеств~$X$ таких, что
$X = \bigcup_{i \in I} U_i$.
\end{definition}

Мы не предполагаем, что элементы покрытия~$U_i$ являются открытыми
или замкнутыми; вся существенная структура имеет чисто комбинаторный
характер. В приложениях к каузальным множествам $X$ является конечным
частично упорядоченным множеством, а $U_i$ --- окрестностями Александрова
или жадными антицепными слоями~\cite{Sorkin:2003bx,Major:2007}.

\begin{definition}[Нерв]
\label{def:nerve}
\emph{Нерв} конечного покрытия~$\cU = (U_i)_{i \in I}$ --- это абстрактный
симплициальный комплекс~$\cN = \cN(\cU)$ на множестве
вершин~$I$, симплексами которого являются непустые конечные
подмножества~$\sigma \subseteq I$ такие, что
$\bigcap_{i \in \sigma} U_i \neq \varnothing$.
\end{definition}

\begin{definition}[Симплексы и грани]
\label{def:simplices}
Симплекс~$\sigma \in \cN$ мощности~$|\sigma| = k + 1$ называется
\emph{$k$-симплексом}. Через~$\cN_k$ обозначается множество $k$-симплексов.
В частности:
\begin{itemize}[nosep]
  \item $\cN_0$ состоит из одноэлементных множеств $\{i\}$, $i \in I$ (вершины);
  \item $\cN_1$ состоит из пар $\{i,j\}$, где $i \neq j$ и
        $U_i \cap U_j \neq \varnothing$ (рёбра);
  \item $\cN_2$ состоит из троек $\{i,j,k\}$ с попарно непустыми
        пересечениями и $U_i \cap U_j \cap U_k \neq \varnothing$ (треугольники).
\end{itemize}
\emph{Грань коразмерности один} $k$-симплекса~$\sigma$ --- это
$(k{-}1)$-симплекс~$\tau \subset \sigma$ с $|\tau| = k$.
\end{definition}

\begin{definition}[Ветвящее ребро]
\label{def:branching}
Ребро~$e \in \cN_1$ называется \emph{треугольно-несущим}, если существует
хотя бы один треугольник~$\sigma \in \cN_2$ такой, что~$e$ является
гранью коразмерности один в~$\sigma$. Ребро называется \emph{ветвящим},
если оно является гранью коразмерности один хотя бы двух различных
треугольников.
\end{definition}

\begin{figure}[t]
\centering
\begin{tikzpicture}[
  vertex/.style={circle, draw, fill=black, inner sep=1.5pt, minimum size=5pt},
  edge/.style={thick},
  branchedge/.style={very thick, blue!70!black},
  triA/.style={fill=blue!8},
  triB/.style={fill=red!8},
  label distance=3pt
]
  % Вершины
  \node[vertex, label=left:$a$]  (a) at (0,0) {};
  \node[vertex, label=right:$b$] (b) at (3.2,0) {};
  \node[vertex, label=above:$c$] (c) at (1.6,2.2) {};
  \node[vertex, label=below:$d$] (d) at (1.6,-2.2) {};
  \node[vertex, label=left:$f$]  (f) at (-1.8,1.2) {};

  % Закрашенные треугольники
  \begin{scope}[on background layer]
    \fill[triA] (a.center) -- (b.center) -- (c.center) -- cycle;
    \fill[triB] (a.center) -- (b.center) -- (d.center) -- cycle;
  \end{scope}

  % Рёбра
  \draw[branchedge] (a) -- (b) node[midway, above, font=\footnotesize,
    text=blue!70!black] {ветвящее};
  \draw[edge] (a) -- (c);
  \draw[edge] (b) -- (c);
  \draw[edge] (a) -- (d);
  \draw[edge] (b) -- (d);
  \draw[edge] (a) -- (f);

  % Подписи треугольников
  \node[font=\small] at (1.6,0.9) {$T_1 = \{a,b,c\}$};
  \node[font=\small] at (1.6,-0.9) {$T_2 = \{a,b,d\}$};

\end{tikzpicture}
\caption{Конечный нерв с пятью вершинами, шестью рёбрами и двумя
треугольниками. Ребро~$\{a,b\}$ является \emph{ветвящим ребром}: оно
является гранью коразмерности один как~$T_1$, так и~$T_2$. Вершина~$f$
соединена только с~$a$; ребро $\{a,f\}$ не является ни треугольно-несущим,
ни ветвящим.}
\label{fig:nerve}
\end{figure}

\begin{example}
\label{ex:nerve}
На рис.~\ref{fig:nerve} показан нерв~$\cN$ с множеством вершин
$I = \{a,b,c,d,f\}$. Нерв имеет два треугольника, $T_1 = \{a,b,c\}$ и
$T_2 = \{a,b,d\}$, которые имеют общее ветвящее ребро~$\{a,b\}$.
Остальные рёбра~$\{a,c\}$, $\{b,c\}$, $\{a,d\}$, $\{b,d\}$ являются
треугольно-несущими, но не ветвящими (каждое принадлежит ровно одному
треугольнику). Ребро~$\{a,f\}$ не является треугольно-несущим.
\end{example}

% ============================================================
\section{Данные ориентации и граничные операторы}
\label{sec:orientation}
% ============================================================

Для определения граничных операторов на абстрактном симплициальном
комплексе, симплексы которого являются \emph{неупорядоченными} конечными
множествами, необходимо ввести вспомогательные \emph{данные ориентации},
задающие упорядочение вершин каждого симплекса. Мы напоминаем стандартную
конструкцию (ср.~\cite[глава~2]{Hatcher:2002},
\cite[глава~1]{Munkres:1984}) и адаптируем её к случаю конечного нерва.

\begin{definition}[Упорядоченный свидетель]
\label{def:witness}
\emph{Упорядоченный свидетель} $k$-симплекса~$\sigma = \{v_0, \ldots, v_k\}
\in \cN_k$ --- это биекция $w \colon \{0, \ldots, k\} \to \sigma$.
Мы записываем~$w = (w(0), w(1), \ldots, w(k))$ как упорядоченный кортеж.
\end{definition}

Для $1$-симплекса~$e = \{a,b\}$ два упорядоченных свидетеля ---
это~$(a,b)$ и~$(b,a)$. Для $2$-симплекса~$\sigma = \{a,b,c\}$ существует
шесть упорядоченных свидетелей, соответствующих шести перестановкам~$\{a,b,c\}$.

\begin{definition}[Индуцированная граница ребра]
\label{def:edge-boundary}
Пусть задан упорядоченный свидетель~$w = (v_0, v_1)$ ребра~$e \in \cN_1$.
\emph{Индуцированная граница} определяется как
\begin{equation}
  \label{eq:edge-boundary}
  \partial_1 [v_0, v_1] \;=\; [v_1] - [v_0]\,.
\end{equation}
В матричной форме граничный оператор~$\partial_1$ является
$|\cN_0| \times |\cN_1|$ матрицей над~$\ZZ$ с элементами
\[
  (\partial_1)_{v,e} =
  \begin{cases}
    +1 & \text{если } v = w(1) \text{ (положительная вершина)}, \\
    -1 & \text{если } v = w(0) \text{ (отрицательная вершина)}, \\
    \;0 & \text{если } v \notin e.
  \end{cases}
\]
\end{definition}

\begin{definition}[Индуцированная граница треугольника]
\label{def:triangle-boundary}
Пусть задан упорядоченный свидетель~$w = (v_0, v_1, v_2)$
треугольника~$\sigma \in \cN_2$. \emph{Индуцированная граница}
определяется как
\begin{equation}
  \label{eq:triangle-boundary}
  \partial_2 [v_0, v_1, v_2] \;=\; [v_1, v_2] - [v_0, v_2] + [v_0, v_1]\,.
\end{equation}
Точнее, каждое граничное ребро~$\sigma$ наследует ориентацию:
ребро, полученное удалением вершины~$v_j$, получает знак~$(-1)^j$.
\end{definition}

\begin{definition}[Позиционное ребро и выбранная грань]
\label{def:chosen-face}
Пусть~$w = (w_0, w_1, w_2)$ --- упорядоченный свидетель
треугольника~$\sigma \in \cN_2$. Для $0 \leq p < q \leq 2$
\emph{$pq$-ребро} свидетеля~$w$ --- это ребро~$\{w_p, w_q\} \in \cN_1$.
Каждое граничное ребро~$e$ симплекса~$\sigma$ является $pq$-ребром
свидетеля~$w$ для единственной пары~$p < q$, поскольку три граничных
ребра суть $\{w_0,w_1\}$, $\{w_0,w_2\}$ и~$\{w_1,w_2\}$.

\emph{Выбранная грань} граничного ребра~$e = \{w_p, w_q\}$ (при $p < q$)
относительно~$w$ --- это грань коразмерности один~$\{w_q\}$ ребра ---
вершина, занимающая более позднюю позицию в кортеже свидетеля. Мы
записываем
\[
  \chosenface_w(e) = \{w_q\}\,.
\]
Эквивалентно, $\chosenface_w(e)$ --- это вершина ребра~$e$ с
положительным коэффициентом в границе~$\partial_1[w_p, w_q] = [w_q] - [w_p]$.
\end{definition}

\begin{example}
\label{ex:chosen}
Для треугольника~$\sigma = \{a,b,c\}$ со свидетелем~$w = (a,b,c)$:
\begin{itemize}[nosep]
  \item $\chosenface_w(\{a,b\}) = \{b\}$
        \quad (из упорядочения $(a,b)$ вторая вершина ---~$b$),
  \item $\chosenface_w(\{a,c\}) = \{c\}$
        \quad (из упорядочения $(a,c)$ вторая вершина ---~$c$),
  \item $\chosenface_w(\{b,c\}) = \{c\}$
        \quad (из упорядочения $(b,c)$ вторая вершина ---~$c$).
\end{itemize}
При транспозиции~$\swapAB(w) = (b,a,c)$ ребро~$\{a,b\}$
наследует упорядочение~$(b,a)$, поэтому выбранная грань становится
$\chosenface_{\swapAB(w)}(\{a,b\}) = \{a\} \neq \{b\}$.
\end{example}

\begin{definition}[Глобальные данные ориентации]
\label{def:global-orientation}
\emph{Глобальные данные ориентации рёбер} --- это функция, сопоставляющая
каждому ребру~$e \in \cN_1$ упорядоченный свидетель ребра~$e$.
\emph{Глобальные свидетельские данные треугольников} --- это
функция~$\cT$, сопоставляющая каждому треугольнику~$\sigma \in \cN_2$
упорядоченный свидетель~$\cT(\sigma)$ треугольника~$\sigma$.
\end{definition}

Если заданы глобальные данные ориентации, то граничные
операторы~$\partial_1$ и~$\partial_2$ суть целочисленные матрицы,
определяемые применением определений~\ref{def:edge-boundary}
и~\ref{def:triangle-boundary} к каждому симплексу с использованием
назначенного ему упорядоченного свидетеля.

% ============================================================
\section{Совместимость и цепной комплекс}
\label{sec:chain-complex}
% ============================================================

Свойство цепного комплекса~$\partial_1 \circ \partial_2 = 0$ не
выполняется для произвольных данных ориентации. Оно требует \emph{условия
совместимости}, связывающего ориентации рёбер со свидетелями
треугольников.

\begin{definition}[Совместимость]
\label{def:compatibility}
Пусть $\omega_1$ --- глобальные данные ориентации рёбер, а $\omega_2$ ---
глобальные свидетельские данные треугольников. Мы говорим, что
$(\omega_1, \omega_2)$ \emph{совместимы}, если для каждой
вершины~$v \in \cN_0$ и каждого треугольника~$\sigma \in \cN_2$
три знаковых вклада от граничных рёбер~$\sigma$ в вершину~$v$
компенсируются:
\begin{equation}
  \label{eq:composition-entry}
  \sum_{e \in \cN_1} (\partial_1)_{v, e}\, (\partial_2)_{e, \sigma}
  \;=\; 0\,.
\end{equation}
Поскольку~$(\partial_2)_{e,\sigma} \neq 0$ лишь тогда, когда~$e$ является
граничным ребром~$\sigma$, а каждый треугольник имеет ровно три граничных
ребра, сумма содержит не более трёх ненулевых слагаемых.
\end{definition}

\begin{theorem}[Цепной комплекс]
\label{thm:chain-complex}
Пусть~$\cN$ --- конечный нерв и~$(\omega_1, \omega_2)$ --- совместимые
глобальные данные ориентации. Тогда:
\begin{enumerate}[label=\textup{(\roman*)},nosep]
  \item Композиция обращается в нуль:~$\partial_1 \circ \partial_2 = 0$
        как матрица над~$\ZZ$.
  \item Каждая $1$-граница является $1$-циклом: $\operatorname{im} \partial_2
        \subseteq \ker \partial_1$, то есть $B_1 \leq Z_1$.
  \item Первая группа гомологий
        \[
          H_1 \;=\; Z_1 / B_1
          \;=\; \ker \partial_1 \,/\, \operatorname{im} \partial_2
        \]
        является корректно определённой конечно порождённой абелевой группой.
\end{enumerate}
\end{theorem}

\begin{proof}[Набросок доказательства]
Часть~(i): элемент $(v, \sigma)$ матрицы~$\partial_1 \circ \partial_2$
есть в точности сумма из~\eqref{eq:composition-entry}, которая обращается
в нуль по условию совместимости. Поскольку это выполняется для
всех~$(v,\sigma)$, вся матрица нулевая. Часть~(ii) следует из~(i): если
$z = \partial_2 c$ для некоторой $2$-цепи~$c$, то
$\partial_1 z = (\partial_1 \circ \partial_2) c = 0$.
Часть~(iii) является стандартным алгебраическим следствием из~(ii).
\end{proof}

Условие совместимости можно эквивалентно выразить в терминах
\emph{локальных данных ориентации}. Локальные данные ориентации
сопоставляют каждому симплексу согласованную разметку его граней
коразмерности один. Ключевая эквивалентность состоит в следующем:

\begin{proposition}[Эквивалентность совместимости]
\label{prop:compat-equiv}
Глобальные данные ориентации рёбер~$\omega_1$ и глобальные свидетельские
данные треугольников~$\omega_2$ совместимы тогда и только тогда, когда
существуют \emph{граничные локальные данные ориентации}~$\theta$ ---
функция, сопоставляющая каждому симплексу согласованную разметку его
граней коразмерности один, --- такие, что~$\omega_1$ и~$\omega_2$
индуцированы из~$\theta$.
\end{proposition}

\begin{proof}[Набросок доказательства]
В прямом направлении~$\theta$ строится из~$(\omega_1, \omega_2)$ путём
считывания знаковых назначений на гранях коразмерности один из граничных
матриц. В обратном направлении проверяется, что если~$\omega_1$
и~$\omega_2$ индуцированы из одного~$\theta$, то для каждой вершины~$v$
и каждого треугольника~$\sigma = [v_0, v_1, v_2]$ сумма
$\sum_{j=0}^{2} (-1)^j (\partial_1)_{v,\, [v_0, \ldots, \hat{v}_j,
\ldots, v_2]}$ обращается в нуль, что является стандартным сокращением
«граница границы».
Полное доказательство содержится в модуле Lean \texttt{BoundaryCompatibilityEquiv}.
\end{proof}

% ============================================================
\section{Существование совместимых данных ориентации}
\label{sec:existence}
% ============================================================

Теперь мы покажем, что совместимые данные ориентации всегда существуют.
Построение проводится в два этапа: сначала каждому треугольнику
назначается упорядоченный свидетель; затем из свидетелей треугольников
выводятся совместимые ориентации рёбер.

\begin{definition}[Совместимость выбора ребро--треугольник]
\label{def:edge-triangle-compat}
Пусть~$\chi$ --- глобальные данные ориентации рёбер, а~$\cT$ ---
глобальные свидетельские данные треугольников. Мы говорим, что
$(\chi, \cT)$ \emph{глобально совместимы по выбору}, если для каждого
треугольника $\sigma \in \cN_2$ и каждого граничного ребра~$e$
треугольника~$\sigma$ ориентация ребра, назначенная в~$\chi$, согласуется
с ориентацией, индуцированной на~$e$ свидетелем
треугольника~$\cT(\sigma)$.
\end{definition}

\begin{construction}[Совместимые данные из свидетелей треугольников]
\label{con:from-witnesses}
Пусть~$\cT$ --- глобальные свидетельские данные треугольников,
удовлетворяющие \emph{условию когерентности} (которое будет определено
в разделе~\ref{sec:coherence}). Для каждого ребра~$e \in \cN_1$:
\begin{itemize}[nosep]
  \item Если~$e$ является граничным ребром некоторого
        треугольника~$\sigma$, определяем ориентацию ребра~$e$ как
        ориентацию, индуцированную свидетелем~$\cT(\sigma)$.
        Условие когерентности гарантирует, что результат не зависит
        от выбора~$\sigma$.
  \item Если~$e$ не является треугольно-несущим, назначаем~$e$
        произвольную ориентацию.
\end{itemize}
Это порождает глобально совместимую по выбору пару~$(\chi, \cT)$.
\end{construction}

\begin{proposition}
\label{prop:coherent-exists}
Для любого конечного нерва~$\cN$ существуют глобальные свидетельские
данные треугольников~$\cT$, удовлетворяющие условию когерентности из
определения~\ref{def:coherence}.
\end{proposition}

\begin{proof}[Набросок доказательства]
Выберем произвольный линейный порядок~$\leq$ на множестве вершин~$I$.
Для каждого треугольника $\sigma = \{i,j,k\}$, где~$i < j < k$,
положим $\cT(\sigma) = (i,j,k)$.
Полученные данные когерентны, поскольку любые два треугольника,
имеющих общее ребро~$\{i,j\}$, оба помещают~$i$ перед~$j$, а значит,
согласуются в выбранной грани~$\{j\}$. Эта конструкция неканонична:
она зависит от выбора линейного порядка на~$I$.
\end{proof}

\begin{theorem}[Существование]
\label{thm:existence}
Для любого конечного нерва~$\cN$ существуют совместимые данные
ориентации $(\omega_1, \omega_2)$ такие, что
$\partial_1 \circ \partial_2 = 0$.
\end{theorem}

\begin{proof}
Возьмём когерентные глобальные свидетельские данные~$\cT$ из
утверждения~\ref{prop:coherent-exists} и применим
конструкцию~\ref{con:from-witnesses} для получения глобально совместимой
по выбору пары~$(\omega_1, \omega_2)$. Когерентность гарантирует,
что ориентация рёбер на каждом треугольно-несущем ребре корректно
определена, а непосредственная проверка~\eqref{eq:composition-entry}
для каждого~$(v, \sigma)$ даёт~$\partial_1 \circ \partial_2 = 0$.
\end{proof}

\begin{remark}
\label{rem:noncanonical}
Ключевой момент состоит в том, что доказательство существования
использует линейный порядок на множестве вершин, который \emph{не является
частью данных нерва}. Структура нерва предоставляет лишь неупорядоченные
симплексы. Вопрос о том, может ли канонический линейный порядок быть
выведен из более богатых структурных данных (таких как частичный порядок
на базовом пространстве), остаётся открытым; см.
раздел~\ref{sec:discussion}.
\end{remark}

% ============================================================
\section{Условия когерентности и их эквивалентности}
\label{sec:coherence}
% ============================================================

Совместимость конструкции~\ref{con:from-witnesses} требует, чтобы все
свидетели треугольников на общем ребре согласовывались в индуцированной
ориентации ребра. Теперь мы сформулируем это точно и установим
несколько эквивалентных формулировок.

\begin{definition}[Попарная ветвящая когерентность]
\label{def:coherence}
Пусть~$\cT$ --- глобальные свидетельские данные треугольников. Мы
говорим, что~$\cT$ удовлетворяет \emph{попарной когерентности ветвящих
выбранных граней}, если для каждого ветвящего ребра~$e \in \cN_1$
и каждой пары треугольников~$\sigma_1, \sigma_2 \in \cN_2$, имеющих~$e$
в качестве грани коразмерности один, выбранные грани совпадают:
\[
  \chosenface_{\cT(\sigma_1)}(e)
  \;=\; \chosenface_{\cT(\sigma_2)}(e)\,.
\]
\end{definition}

\begin{definition}[Конфликт на ветвящем ребре]
\label{def:conflict}
\emph{Конфликт на ветвящем ребре} для глобальных свидетельских
данных~$\cT$ --- это ветвящее ребро~$e$ и два треугольника~$\sigma_1,
\sigma_2$, имеющих~$e$ в качестве грани, такие, что индуцированные
локальные ориентации ребра не согласуются:
\[
  \chosenface_{\cT(\sigma_1)}(e)
  \;\neq\; \chosenface_{\cT(\sigma_2)}(e)\,.
\]
\end{definition}

\begin{proposition}[Двойственность конфликта и когерентности]
\label{prop:conflict-iff}
Глобальные свидетельские данные~$\cT$ удовлетворяют попарной ветвящей
когерентности тогда и только тогда, когда они не содержат конфликтов
на ветвящих рёбрах.
\end{proposition}

\begin{proof}
Непосредственно следует из определений: когерентность есть отрицание
конфликта на каждом ветвящем ребре.
\end{proof}

Условие когерентности допускает несколько эквивалентных формулировок,
каждая из которых освещает различные аспекты:

\begin{theorem}[Эквивалентность условий когерентности]
\label{thm:equiv}
Пусть~$\cT$ --- глобальные свидетельские данные треугольников. Следующие
условия эквивалентны:
\begin{enumerate}[label=\textup{(\roman*)},nosep]
  \item \emph{Попарная ветвящая когерентность:}
        все свидетели треугольников на каждом ветвящем ребре индуцируют
        одну и ту же выбранную грань.
  \item \emph{Склеенные выборы ветвящих рёбер:}
        существует функция, сопоставляющая каждому ветвящему ребру
        выбранную грань, согласованная со всеми примыкающими свидетелями
        треугольников.
  \item \emph{Чисто треугольная когерентность перекрытий:}
        для каждой пары различных треугольников, имеющих общую грань
        коразмерности один, индуцированные локальные ориентации на
        общем ребре совпадают.
  \item \emph{Глобальная когерентность ориентации рёбер:}
        свидетели треугольников индуцируют корректно определённую
        глобальную ориентацию на каждом треугольно-несущем ребре.
\end{enumerate}
\end{theorem}

\begin{proof}[Набросок доказательства]
Импликации (ii)$\Rightarrow$(i) и (i)$\Rightarrow$(iii)
получаются раскрытием определений.
Импликация (iii)$\Rightarrow$(ii) использует тот факт, что на каждом
ветвящем ребре попарное согласование всех свидетелей треугольников
влечёт существование единой согласованной выбранной грани.
Эквивалентность (i)$\Leftrightarrow$(iv) следует из наблюдения, что
локальная ориентация ребра определяется выбранной гранью
(через~\eqref{eq:edge-boundary}), и обратно, выбранная грань
восстанавливается из ориентации. Полные доказательства содержатся
в формализации на Lean; см. приложение~\ref{app:lean}.
\end{proof}

\begin{corollary}
\label{cor:coherence-implies-compat}
Если глобальные свидетельские данные~$\cT$ удовлетворяют попарной
ветвящей когерентности, то существуют совместимые данные ориентации
$(\omega_1, \omega_2)$ с~$\partial_1 \circ \partial_2 = 0$, где
$\omega_2$ определяется данными~$\cT$, а $\omega_1$ --- индуцированная
глобальная ориентация рёбер.
\end{corollary}

\begin{proof}
Применим конструкцию~\ref{con:from-witnesses} для получения~$\omega_1$
из~$\cT$: условие когерентности (определение~\ref{def:coherence})
гарантирует корректную определённость индуцированной ориентации рёбер
на каждом треугольно-несущем ребре. Полученная
пара~$(\omega_1, \omega_2)$ глобально совместима по выбору
(определение~\ref{def:edge-triangle-compat}), что влечёт совместимость
(определение~\ref{def:compatibility}).
\end{proof}

\begin{remark}[Один нерв, разные упорядочения]
\label{rem:same-support}
Поскольку глобальные свидетельские данные --- это функция, определённая
на всём~$\cN_2$, любые два набора глобальных свидетельских данных на
одном и том же нерве~$\cN$ автоматически определены на одной и той же
области треугольников.
Они могут различаться лишь \emph{упорядочениями вершин}, назначенными
отдельным треугольникам, но не тем, какие треугольники получают свидетелей.
Когда мы говорим, что~$\cT$ и~$\cT'$ «имеют один и тот же симплициальный
носитель», мы имеем в виду именно это: они являются глобальными
свидетельскими данными на одном и том же нерве.
\end{remark}

% ============================================================
\section{Препятствие невнутренности}
\label{sec:obstruction}
% ============================================================

Теперь мы докажем основной отрицательный результат: условие когерентности
из раздела~\ref{sec:coherence} не является свойством симплициальных данных
носителя нерва. Оно зависит от конкретных упорядоченных свидетелей,
назначенных каждому треугольнику, и различные выборы свидетелей на одном
и том же симплициальном носителе могут приводить к взаимоисключающим
результатам когерентности.

\subsection{Транспозиции вершин свидетелей треугольников}
\label{sec:swaps}

\begin{definition}[Транспозиции вершин]
\label{def:swaps}
Пусть~$w = (v_0, v_1, v_2)$ --- упорядоченный свидетель
треугольника~$\sigma$. Определим три операции транспозиции:
\begin{align}
  \swapAB(w) &= (v_1, v_0, v_2)\,, \label{eq:swapAB} \\
  \swapAC(w) &= (v_2, v_1, v_0)\,, \label{eq:swapAC} \\
  \swapBC(w) &= (v_0, v_2, v_1)\,. \label{eq:swapBC}
\end{align}
Каждая транспозиция порождает допустимый упорядоченный свидетель того же
базового треугольника~$\sigma$.
\end{definition}

Ключевое наблюдение состоит в том, что транспозиции сохраняют базовое
ребро, но изменяют выбранную грань:

\begin{lemma}[Сохранение ребра]
\label{lem:edge-preservation}
Пусть~$w = (v_0, v_1, v_2)$. Тогда:
\begin{enumerate}[label=\textup{(\alph*)},nosep]
  \item $\swapAB(w)$ сохраняет $01$-ребро
        (определение~\ref{def:chosen-face}): как~$w$, так и~$\swapAB(w)$
        имеют~$\{v_0, v_1\}$ в качестве $01$-ребра.
  \item $\swapAC(w)$ сохраняет $02$-ребро $\{v_0, v_2\}$.
  \item $\swapBC(w)$ сохраняет $12$-ребро $\{v_1, v_2\}$.
\end{enumerate}
\end{lemma}

\begin{proof}
В каждом случае транспозиция меняет местами две вершины ребра, что
сохраняет ребро как неупорядоченное множество.
\end{proof}

\begin{theorem}[Гибкость свидетеля]
\label{thm:flexibility}
Пусть~$w = (v_0, v_1, v_2)$ --- упорядоченный свидетель. Тогда
справедливы следующие утверждения:
\begin{enumerate}[label=\textup{(\alph*)},nosep]
  \item $\chosenface_w(\{v_0, v_1\}) \neq
        \chosenface_{\swapAB(w)}(\{v_0, v_1\})$.
  \item $\chosenface_w(\{v_0, v_2\}) \neq
        \chosenface_{\swapAC(w)}(\{v_0, v_2\})$.
  \item $\chosenface_w(\{v_1, v_2\}) \neq
        \chosenface_{\swapBC(w)}(\{v_1, v_2\})$.
\end{enumerate}
Иными словами, каждая транспозиция меняет выбранную грань на сохранённом
ребре на противоположную.
\end{theorem}

\begin{proof}
Докажем (a); рассуждения для (b) и (c) аналогичны.
При свидетеле~$w = (v_0, v_1, v_2)$ ребро~$\{v_0, v_1\}$ наследует
упорядочение~$(v_0, v_1)$, поэтому
$\chosenface_w(\{v_0, v_1\}) = \{v_1\}$.
При~$\swapAB(w) = (v_1, v_0, v_2)$ то же ребро наследует
упорядочение~$(v_1, v_0)$, поэтому
$\chosenface_{\swapAB(w)}(\{v_0, v_1\}) = \{v_0\}$.
Поскольку~$v_0 \neq v_1$ (вершины ребра различны),
$\{v_1\} \neq \{v_0\}$.
\end{proof}

\subsection{Основная теорема о препятствии}
\label{sec:main-theorem}

Теперь мы поднимаем гибкость свидетеля одного треугольника до глобального
утверждения о парах треугольников с общим ветвящим ребром.

\begin{definition}[Замена свидетеля]
\label{def:replacement}
Пусть~$\cT$ --- глобальные свидетельские данные треугольников. Для
треугольника~$\sigma$ и упорядоченного свидетеля~$w$
треугольника~$\sigma$ обозначим через~$\cT[\sigma \leftarrow w]$
глобальные данные, совпадающие с~$\cT$ на всех треугольниках,
кроме~$\sigma$, где они назначают~$w$.
\end{definition}

\begin{theorem}[Препятствие невнутренности]
\label{thm:main-obstruction}
Пусть~$\cN$ --- конечный нерв. Пусть~$T_1 \neq T_2 \in \cN_2$ ---
различные треугольники с общим граничным ребром~$e$. Пусть~$w_1$ ---
упорядоченный свидетель~$T_1$, а~$w_2$ --- упорядоченный свидетель~$T_2$,
причём~$e$ является $01$-ребром как~$w_1$, так и~$w_2$. Для любых
базовых глобальных свидетельских данных~$\cT_0$ определим:
\begin{align}
  \cT &\;=\; \cT_0[T_1 \leftarrow w_1][T_2 \leftarrow w_2]\,,
  \label{eq:tau} \\
  \cT' &\;=\; \cT_0[T_1 \leftarrow w_1][T_2 \leftarrow \swapAB(w_2)]\,.
  \label{eq:tau-swap}
\end{align}
Тогда:
\begin{enumerate}[label=\textup{(\alph*)},nosep]
  \item\label{item:not-both}
    $\cT$ и~$\cT'$ не могут одновременно удовлетворять попарной ветвящей
    когерентности на ребре~$e$:
    \[
      \neg\bigl(\text{$\cT$ когерентно на $e$}
      \;\wedge\;
      \text{$\cT'$ когерентно на $e$}\bigr)\,.
    \]
  \item\label{item:conditional}
    Если~$\cT$ удовлетворяет попарной ветвящей когерентности на~$e$,
    то~$\cT'$ содержит конфликт на ветвящем ребре~$e$.
  \item\label{item:global}
    $\cT$ и~$\cT'$ не могут одновременно удовлетворять \emph{глобальной}
    попарной ветвящей когерентности:
    \[
      \neg\bigl(\text{$\cT$ глобально когерентно}
      \;\wedge\;
      \text{$\cT'$ глобально когерентно}\bigr)\,.
    \]
\end{enumerate}
\end{theorem}

\begin{proof}
Поскольку~$T_1 \neq T_2$ и оба имеют~$e$ в качестве грани коразмерности
один, ребро~$e$ является ветвящим.
Запишем~$w_2 = (w_2^{(0)}, w_2^{(1)}, w_2^{(2)})$. По условию,
$e = \{w_2^{(0)}, w_2^{(1)}\}$ является $01$-ребром свидетеля~$w_2$.

\medskip\noindent\textit{Ключевое наблюдение.}
Согласно определению~\ref{def:chosen-face}:
\begin{align*}
  \chosenface_{w_2}(e) &= \{w_2^{(1)}\}\,,\\
  \chosenface_{\swapAB(w_2)}(e) &= \{w_2^{(0)}\}\,.
\end{align*}
Второе равенство выполняется, поскольку~$\swapAB(w_2)
= (w_2^{(1)}, w_2^{(0)}, w_2^{(2)})$ имеет то же базовое ребро
$\{w_2^{(1)}, w_2^{(0)}\} = e$, но с переставленными позициями, так что
вершина на более поздней позиции теперь~$w_2^{(0)}$.
Поскольку~$w_2^{(0)} \neq w_2^{(1)}$ (вершины ребра различны):
\begin{equation}
  \label{eq:two-faces-differ}
  \chosenface_{w_2}(e) \;\neq\; \chosenface_{\swapAB(w_2)}(e)\,.
\end{equation}

\medskip\noindent\textit{Доказательство~\ref{item:not-both}.}
Для данных~$\cT$ свидетели на~$T_1$ и~$T_2$ индуцируют выбранные грани
$\chosenface_{w_1}(e)$ и $\chosenface_{w_2}(e)$ на общем ребре~$e$.
Для данных~$\cT'$ они индуцируют $\chosenface_{w_1}(e)$ и
$\chosenface_{\swapAB(w_2)}(e)$.
Когерентность на~$e$ для~$\cT$ требует
$\chosenface_{w_1}(e) = \chosenface_{w_2}(e)$;
когерентность на~$e$ для~$\cT'$ требует
$\chosenface_{w_1}(e) = \chosenface_{\swapAB(w_2)}(e)$.
Из~\eqref{eq:two-faces-differ} правые части различны,
поэтому~$\chosenface_{w_1}(e)$ не может быть равна обеим одновременно.

\medskip\noindent\textit{Доказательство~\ref{item:conditional}.}
Если~$\cT$ когерентно на~$e$, то
$\chosenface_{w_1}(e) = \chosenface_{w_2}(e) \neq
\chosenface_{\swapAB(w_2)}(e)$, следовательно,~$\cT'$ содержит конфликт
на~$e$ (определение~\ref{def:conflict}).

\medskip\noindent\textit{Доказательство~\ref{item:global}.}
Глобальная когерентность влечёт локальную когерентность на каждом ветвящем
ребре, в том числе на~$e$, поэтому~\ref{item:global}
следует из~\ref{item:not-both}.
\end{proof}

\begin{figure}[t]
\centering
\begin{tikzpicture}[
  vertex/.style={circle, draw, fill=black, inner sep=1.5pt, minimum size=5pt},
  posvertex/.style={circle, draw=green!60!black, fill=green!30,
    inner sep=1.5pt, minimum size=7pt},
  negvertex/.style={circle, draw=red!60!black, fill=red!20,
    inner sep=1.5pt, minimum size=7pt},
  edge/.style={thick},
  branchedge/.style={very thick, blue!70!black},
  checkmark/.style={font=\Large},
  panel/.style={font=\bfseries\small}
]
  % === Панель (a): Когерентно ===
  \begin{scope}[shift={(0,0)}]
    \node[panel] at (1.5,3.2) {(a) Когерентно};

    % Вершины
    \node[posvertex, label={[label distance=1pt]right:\small$b$}]
      (a-b) at (3,0) {};
    \node[negvertex, label={[label distance=1pt]left:\small$a$}]
      (a-a) at (0,0) {};
    \node[vertex, label=above:\small$c$] (a-c) at (1.5,2.2) {};
    \node[vertex, label=below:\small$d$] (a-d) at (1.5,-2.2) {};

    % Треугольники
    \fill[blue!6] (a-a.center) -- (a-b.center) -- (a-c.center) -- cycle;
    \fill[red!6]  (a-a.center) -- (a-b.center) -- (a-d.center) -- cycle;

    % Рёбра
    \draw[branchedge] (a-a) -- (a-b);
    \draw[edge] (a-a) -- (a-c); \draw[edge] (a-b) -- (a-c);
    \draw[edge] (a-a) -- (a-d); \draw[edge] (a-b) -- (a-d);

    % Стрелки свидетелей
    \draw[-{Stealth[length=4pt]}, thick, blue!50!black]
      ($(a-a)+(0.15,0.15)$) -- ($(a-b)+(-0.15,0.15)$);
    \draw[-{Stealth[length=4pt]}, thick, red!50!black]
      ($(a-a)+(0.15,-0.15)$) -- ($(a-b)+(-0.15,-0.15)$);

    % Подписи
    \node[font=\scriptsize, blue!50!black] at (1.5, 1.0)
      {$w_1\!=\!(a,b,c)$};
    \node[font=\scriptsize, red!50!black] at (1.5, -1.0)
      {$w_2\!=\!(a,b,d)$};
    \node[font=\scriptsize] at (1.5, 0.35)
      {оба выбирают $\{b\}$};

    % Галочка
    \node[checkmark, green!50!black] at (1.5, -3.1) {\checkmark};
  \end{scope}

  % === Панель (b): Конфликт ===
  \begin{scope}[shift={(7,0)}]
    \node[panel] at (1.5,3.2) {(b) Конфликт};

    % Вершины — обе выделены как спорные
    \node[vertex, draw=orange!80!black, fill=orange!20,
      inner sep=1.5pt, minimum size=7pt,
      label={[label distance=1pt]right:\small$b$}]
      (b-b) at (3,0) {};
    \node[vertex, draw=orange!80!black, fill=orange!20,
      inner sep=1.5pt, minimum size=7pt,
      label={[label distance=1pt]left:\small$a$}]
      (b-a) at (0,0) {};
    \node[vertex, label=above:\small$c$] (b-c) at (1.5,2.2) {};
    \node[vertex, label=below:\small$d$] (b-d) at (1.5,-2.2) {};

    % Треугольники
    \fill[blue!6] (b-a.center) -- (b-b.center) -- (b-c.center) -- cycle;
    \fill[red!6]  (b-a.center) -- (b-b.center) -- (b-d.center) -- cycle;

    % Рёбра
    \draw[branchedge] (b-a) -- (b-b);
    \draw[edge] (b-a) -- (b-c); \draw[edge] (b-b) -- (b-c);
    \draw[edge] (b-a) -- (b-d); \draw[edge] (b-b) -- (b-d);

    % Стрелки свидетелей (T1 — без изменений, T2 — развёрнут)
    \draw[-{Stealth[length=4pt]}, thick, blue!50!black]
      ($(b-a)+(0.15,0.15)$) -- ($(b-b)+(-0.15,0.15)$);
    \draw[-{Stealth[length=4pt]}, thick, red!50!black]
      ($(b-b)+(-0.15,-0.15)$) -- ($(b-a)+(0.15,-0.15)$);

    % Подписи
    \node[font=\scriptsize, blue!50!black] at (1.5, 1.0)
      {$w_1\!=\!(a,b,c)$};
    \node[font=\scriptsize, red!50!black] at (1.5, -1.0)
      {$w_2'\!=\!(b,a,d)$};

    % Аннотация конфликта: каждая подпись рядом с ВЫБРАННОЙ вершиной
    \node[font=\scriptsize, blue!50!black] at (3.3, 0.45)
      {$w_1 \to \{b\}$};
    \node[font=\scriptsize, red!50!black] at (-0.3, -0.45)
      {$w_2' \to \{a\}$};
    \node[font=\small, red!70!black] at (1.5, 0.38)
      {$\neq$};

    % Крестик
    \node[checkmark, red!70!black] at (1.5, -3.1) {$\times$};
  \end{scope}
\end{tikzpicture}
\caption{Препятствие при замене свидетеля на общем ветвящем ребре
$e = \{a,b\}$. \textbf{(a)}~Оба свидетеля треугольников~$(a,b,c)$
и~$(a,b,d)$ индуцируют одну и ту же выбранную грань~$\{b\}$
(зелёная вершина) на~$e$: данные когерентны. \textbf{(b)}~После
применения~$\swapAB$ к свидетелю~$T_2$ свидетель становится~$(b,a,d)$,
и выбранная грань на~$e$ меняется на~$\{a\}$. Поскольку~$w_1$
по-прежнему выбирает~$\{b\}$, данные содержат конфликт. По
теореме~\ref{thm:main-obstruction}, $\cT$~и~$\cT'$ не могут быть
одновременно когерентными.}
\label{fig:obstruction}
\end{figure}

\begin{corollary}[Невнутренность]
\label{cor:non-intrinsic}
Условие глобальной когерентности (определение~\ref{def:coherence}) не
определяется симплициальным носителем нерва. А именно, существуют
глобальные свидетельские данные~$\cT$ и~$\cT'$ с одним и тем же
симплициальным носителем (одним и тем же множеством треугольников как
неупорядоченных множеств), но с различным статусом когерентности.
\end{corollary}

\begin{proof}
Возьмём любой нерв с ветвящим ребром, выберем
треугольники~$T_1 \neq T_2$, имеющие общее ветвящее ребро, и построим
$\cT$ и~$\cT'$ как в теореме~\ref{thm:main-obstruction}. Данные~$\cT$
и~$\cT'$ совпадают на каждом треугольнике, кроме~$T_2$, где они
назначают различные упорядоченные свидетели \emph{одного и того же
неупорядоченного треугольника}. Таким образом, симплициальный носитель
совпадает. По теореме~\ref{thm:main-obstruction} не более чем одно
из данных~$\cT$ и~$\cT'$ может быть когерентным.
\end{proof}

\begin{remark}
\label{rem:not-impossibility}
Теорема~\ref{thm:main-obstruction} \emph{не} утверждает, что когерентные
глобальные свидетельские данные не существуют. Она утверждает, что
когерентность не является свойством симплициального носителя: среди
свидетельских данных с одним и тем же носителем когерентные образуют
собственное подмножество, зависящее от выбора упорядочений.
Утверждение~\ref{prop:coherent-exists} подтверждает, что когерентные
данные существуют (при использовании линейного порядка на вершинах).
\end{remark}

\begin{remark}
\label{rem:all-edges}
Препятствие однородно по всем трём граничным рёбрам.
Теорема~\ref{thm:flexibility} показывает, что~$\swapAB$, $\swapAC$
и~$\swapBC$ меняют выбранную грань на соответствующих сохранённых рёбрах
на противоположную. Таким образом, невнутренность не является артефактом
конкретной нумерации рёбер.
\end{remark}

\begin{remark}[Локализация]
\label{rem:localization}
Части~\ref{item:conditional} и~\ref{item:global}
теоремы~\ref{thm:main-obstruction} показывают, что препятствие
\emph{одностороннее и локализованное}: если одно из двух свидетельских
данных глобально когерентно, то другое содержит конфликт на том самом
ветвящем ребре, на котором была произведена транспозиция.
\end{remark}

% ============================================================
\section{Обсуждение}
\label{sec:discussion}
% ============================================================

\subsection{Сводка результатов}

Мы установили два дополняющих друг друга результата о граничных операторах
на конечных нервах:

\begin{enumerate}[nosep]
  \item \emph{Вспомогательная конструкция работает.}
    При наличии глобальных свидетельских данных треугольников,
    удовлетворяющих попарной ветвящей когерентности, индуцированные
    граничные операторы образуют цепной комплекс
    с~$\partial_1 \circ \partial_2 = 0$ и корректно определённой первой
    гомологией~$H_1$ (теорема~\ref{thm:chain-complex}).
    Совместимые данные всегда существуют (теорема~\ref{thm:existence}),
    но их построение неканонично.

  \item \emph{Внутренняя когерентность из одних данных носителя
    невозможна.}
    Условие когерентности, обеспечивающее совместимость, не определяется
    симплициальным носителем нерва
    (следствие~\ref{cor:non-intrinsic}).
    Два свидетельских данных на одном и том же носителе, различающихся
    единственной транспозицией вершин, не могут быть одновременно
    когерентными (теорема~\ref{thm:main-obstruction}).
    Это препятствие локализовано на ветвящих рёбрах, однородно по
    всем трём граничным рёбрам и допускает как дизъюнктивную, так и
    условную формы.
\end{enumerate}

\subsection{Следствия для программы спектрального действия}

В контексте спектральной каузальной теории вопрос, мотивирующий данную
работу, состоит в том, достаточно ли комбинаторных данных конечного
приближения пространства-времени для внутреннего определения граничных
операторов и гомологических инвариантов. Наши результаты показывают, что
\emph{одной структуры нерва недостаточно}: неупорядоченный симплициальный
носитель не определяет необходимые данные ориентации.

Это не закрывает более широкую проблему. Три направления остаются
открытыми:

\begin{enumerate}[nosep]
  \item \emph{Выход через каузальный порядок.}
    Конечное каузальное множество несёт частичный порядок, который
    строго богаче данных носителя нерва. Этот частичный порядок
    мог бы в принципе предоставить канонический источник упорядочений
    вершин (например, путём продолжения частичного порядка до линейного
    и использования его для ориентирования всех симплексов). Формализация
    на Lean работает с абстрактными покрытиями и нервами без ссылок на
    частичные порядки, поэтому теорема~\ref{thm:main-obstruction} не
    затрагивает эту возможность.

  \item \emph{Более слабые цели реконструкции.}
    Граничные операторы над~$\ZZ/2\ZZ$ не требуют ориентаций:
    беззнаковые матрицы инцидентности уже определяют цепной комплекс
    над~$\ZZ/2\ZZ$. Результирующая гомология по модулю~$2$ является
    внутренней для носителя нерва. Однако целочисленная гомология
    необходима для большинства физических приложений (ориентация,
    характеристические классы, функционал спектрального действия).

  \item \emph{Ансамблевые и эмерджентные подходы.}
    Можно ослабить требование от внутренней реконструкции на одном нерве
    до статистического проявления из ансамбля нервов или допустить
    определение спектральной тройки лишь после шага огрубления.
    Эти подходы ослабляют непосредственность реконструкции, но
    могут обойти препятствие.
\end{enumerate}

\subsection{Связь с классической теорией ориентации}

Невнутренность данных ориентации на симплициальных комплексах ---
хорошо известное явление в алгебраической топологии
(ср.~\cite{Hatcher:2002,Munkres:1984}). Любой абстрактный симплициальный
комплекс допускает ориентацию (путём выбора линейного порядка на
вершинах), но ориентация не является частью комбинаторных данных. Два
ключевых различия между классической постановкой и рассматриваемой здесь
состоят в следующем:

\begin{itemize}[nosep]
  \item В классической симплициальной гомологии группы гомологий
        не зависят от выбора ориентации (с точностью до изоморфизма).
        В рассматриваемой постановке когерентные данные всегда
        существуют (утверждение~\ref{prop:coherent-exists}), но
        \emph{являются ли конкретные свидетельские данные когерентными}
        зависит от конкретного упорядочения вершин, и
        теорема~\ref{thm:main-obstruction} количественно характеризует
        эту чувствительность.

  \item В классической постановке линейный порядок на вершинах
        даёт каноническую ориентацию. Суть наших результатов в том, что
        в контексте спектральной геометрии спрашивается, может ли такой
        порядок быть выведен из более примитивных данных (каузального
        множества, спектральной тройки), а не наложен извне.
\end{itemize}

\subsection{Открытые вопросы}

В заключение мы формулируем ряд открытых проблем, порождённых данной
работой.

\begin{enumerate}[nosep]
  \item Может ли частичный порядок конечного каузального множества
        канонически определить когерентные глобальные свидетельские данные
        на нерве естественного покрытия?
  \item В стандартной симплициальной гомологии группы гомологий
        не зависят от выбора упорядочения вершин с точностью до
        канонического изоморфизма~\cite[раздел~2.1]{Hatcher:2002}.
        Распространяется ли эта независимость на рассматриваемую
        постановку, где совместимые данные строятся из когерентных
        свидетельских данных посредством конструкции~\ref{con:from-witnesses},
        а не из глобального линейного порядка на вершинах?
  \item Может ли препятствие теоремы~\ref{thm:main-obstruction}
        быть обобщено на симплексы более высокой размерности и
        старшие граничные операторы~$\partial_k$ при~$k \geq 3$?
  \item Существует ли естественное ослабление условия когерентности,
        которое является внутренним для носителя нерва и при этом
        порождает полезный гомологический инвариант?
\end{enumerate}

% ============================================================
\subsection*{Благодарности}

% ============================================================
\appendix

\section{Формализация в Lean 4}
\label{app:lean}

Все определения, формулировки лемм и доказательства теорем данной работы
формализованы в Lean~4 (версия~4.28.0) с использованием библиотеки
Mathlib4. Формализация компилируется командой \texttt{lake build}
и не содержит объявлений~\texttt{sorry}, \texttt{admit}
или~\texttt{axiom} в коде уровня теорем. Исходный код доступен по
адресу \url{https://github.com/davidichalfyorov-wq/sct-theory}.

Формализация состоит из 63 модулей в каталоге
\texttt{SCTLean/FND1/}, организованных следующим образом:

\medskip
\begin{center}
\small
\begin{tabular}{ll}
\toprule
\textbf{Модули} & \textbf{Содержание} \\
\midrule
\texttt{FiniteNerve}, \texttt{SimplicialNerve}, &
  Конструкция нерва, \\
\texttt{TriangleLayer} &
  перечисление симплексов \\[3pt]
\texttt{BoundarySupport},
  \texttt{BoundaryIncidence}, &
  Беззнаковый носитель и \\
\texttt{BoundaryOrientation},
  \texttt{BoundarySignedIncidence} &
  знаковые данные инцидентности \\[3pt]
\texttt{BoundaryTable},
  \texttt{BoundaryMatrix}, &
  Граничные матрицы, \\
\texttt{BoundaryComposition} &
  формула композиции \\[3pt]
\texttt{BoundaryCompatibility},
  \texttt{BoundaryLocalOrientation}, &
  Совместимость и \\
\texttt{BoundaryCompatibleExistence},
  \texttt{BoundaryCompatibilityEquiv} &
  результаты существования \\[3pt]
\texttt{BoundaryEndpointChoice},
  \texttt{BoundaryOrderedTriangle}, &
  Свидетельские данные и \\
\texttt{BoundaryTriangleChoice},
  \texttt{BoundaryChoiceCompatibility} &
  совместимость выборов \\[3pt]
\texttt{BoundaryGlobalChoiceCompatibility}, &
  Глобальная совместимость \\
\texttt{BoundaryGlobalCompositionZero} &
  и $\partial_1 \circ \partial_2 = 0$ \\[3pt]
\texttt{BoundaryChainMaps},
  \texttt{BoundaryHomologyPrelude}, &
  Цепные отображения, \\
\texttt{BoundaryHomologyStructures},
  \texttt{BoundaryHomologyH1} &
  $Z_1$, $B_1$, $H_1$ \\[3pt]
\texttt{BoundaryHomologyClassEq},
  \texttt{BoundaryHomologyUse}, &
  Равенство классов, \\
\texttt{BoundaryHomologyRelabel} &
  перенос при перемаркировке \\[3pt]
\texttt{BoundaryTriangleCoherence}--
  \texttt{BoundaryTriangleOverlapCoherence}, &
  Условия когерентности \\
\texttt{BoundaryPureTriangleOverlapCoherence},
  \texttt{BoundaryBranchingEdge*} &
  и эквивалентности \\[3pt]
\texttt{BoundaryTriangleWitnessFlexibility}, &
  Операции перестановки \\
\texttt{BoundarySharedEdgeWitnessFlexibility}, &
  и основная \\
\texttt{BoundarySharedEdgeTauVariants} &
  теорема о препятствии \\
\bottomrule
\end{tabular}
\end{center}

\medskip

Избранные ключевые имена теорем из формализации с указанием их
соответствия теоремам данной работы:

\medskip
\begin{center}
\small
\begin{tabular}{ll}
\toprule
\textbf{Имя теоремы в Lean} & \textbf{Ссылка в работе} \\
\midrule
\texttt{HasBoundarySquareZero} &
  Теорема~\ref{thm:chain-complex}\,(i) \\
\texttt{oneBoundaries\_le\_oneCycles\_of\_matrix\_zero} &
  Теорема~\ref{thm:chain-complex}\,(ii) \\
\texttt{FirstHomology} &
  Теорема~\ref{thm:chain-complex}\,(iii) \\
\texttt{hasCompatibleLocalOrientation\_of\_} &
  Теорема~\ref{thm:existence} \\
\quad\texttt{globalChoiceCompatible} & \\
\texttt{pureTriangleOverlapCoherence\_iff\_} &
  Теорема~\ref{thm:equiv} \\
\quad\texttt{pairwiseBranchingChosenFaceCoherence} & \\
\texttt{chosenFace\_edgeAB\_differs\_under\_swapAB} &
  Теорема~\ref{thm:flexibility}\,(a) \\
\texttt{not\_both\_pairwiseBranching} &
  Теорема~\ref{thm:main-obstruction} \\
\quad\texttt{ChosenFaceCoherenceAt\_of\_sharedEdge\_swap} & \\
\texttt{branchingEdgeConflictAt\_tauSwap\_of\_} &
  Теорема~\ref{thm:main-obstruction}\,(b) \\
\quad\texttt{pairwiseBranchingChosenFaceCoherenceAt\_} & \\
\quad\texttt{of\_sharedEdge\_swap} & \\
\bottomrule
\end{tabular}
\end{center}

% ============================================================
\begin{thebibliography}{99}

\bibitem{Alfyorov:2026ff}
D.~Alfyorov,
\emph{Nonlocal one-loop form factors of the spectral action with Standard
Model content},
preprint (2026),
\href{https://doi.org/10.22541/au.177507193.33027854/v1}{doi:10.22541/au.177507193.33027854/v1}.

\bibitem{Alfyorov:2026ss}
D.~Alfyorov,
\emph{Solar system and laboratory tests of the spectral action scale},
preprint (2026),
\href{https://doi.org/10.22541/au.177524450.03515205/v1}{doi:10.22541/au.177524450.03515205/v1}.

\bibitem{Alfyorov:2026ct}
D.~Alfyorov,
\emph{Chirality of the Seeley-DeWitt coefficients and quartic Weyl structure
in the spectral action},
preprint (2026),
\href{https://doi.org/10.22541/au.177516440.00576784/v1}{doi:10.22541/au.177516440.00576784/v1}.

\bibitem{Bombelli:1987aa}
L.~Bombelli, J.~Lee, D.~Meyer, and R.~D.~Sorkin,
\emph{Space-time as a causal set},
Phys.\ Rev.\ Lett.\ \textbf{59}, 521 (1987),
\href{https://doi.org/10.1103/PhysRevLett.59.521}{doi:10.1103/PhysRevLett.59.521}.

\bibitem{Chamseddine:1996zu}
A.~H.~Chamseddine and A.~Connes,
\emph{The spectral action principle},
Commun.\ Math.\ Phys.\ \textbf{186}, 731 (1997),
\href{https://doi.org/10.1007/s002200050126}{doi:10.1007/s002200050126},
\href{https://arxiv.org/abs/hep-th/9606001}{arXiv:hep-th/9606001}.

\bibitem{Chamseddine:2006ep}
A.~H.~Chamseddine, A.~Connes, and M.~Marcolli,
\emph{Gravity and the standard model with neutrino mixing},
Adv.\ Theor.\ Math.\ Phys.\ \textbf{11}, 991 (2007),
\href{https://doi.org/10.4310/ATMP.2007.v11.n6.a3}{doi:10.4310/ATMP.2007.v11.n6.a3},
\href{https://arxiv.org/abs/hep-th/0610241}{arXiv:hep-th/0610241}.

\bibitem{Chamseddine:2010ud}
A.~H.~Chamseddine and A.~Connes,
\emph{Noncommutative geometry as a framework for unification of all
fundamental interactions including gravity. Part I.},
Fortschr.\ Phys.\ \textbf{58}, 553 (2010),
\href{https://doi.org/10.1002/prop.201000069}{doi:10.1002/prop.201000069},
\href{https://arxiv.org/abs/1004.0464}{arXiv:1004.0464}.

\bibitem{Connes:1994yd}
A.~Connes,
\emph{Noncommutative Geometry},
Academic Press, San Diego (1994).

\bibitem{Hatcher:2002}
A.~Hatcher,
\emph{Algebraic Topology},
Cambridge University Press (2002).
Available at \url{https://pi.math.cornell.edu/~hatcher/AT/ATpage.html}.

\bibitem{Major:2007}
S.~Major, D.~Rideout, and S.~Surya,
\emph{On recovering continuum topology from a causal set},
J.\ Math.\ Phys.\ \textbf{48}, 032501 (2007),
\href{https://doi.org/10.1063/1.2435599}{doi:10.1063/1.2435599},
\href{https://arxiv.org/abs/gr-qc/0604124}{arXiv:gr-qc/0604124}.

\bibitem{Mathlib2024}
The Mathlib Community,
\emph{Mathlib4},
\url{https://github.com/leanprover-community/mathlib4} (2024).

\bibitem{Munkres:1984}
J.~R.~Munkres,
\emph{Elements of Algebraic Topology},
Addison-Wesley, Menlo Park (1984).

\bibitem{Sorkin:2003bx}
R.~D.~Sorkin,
\emph{Causal sets: discrete gravity},
in \emph{Lectures on Quantum Gravity}, ed.\ A.~Gomberoff and
D.~Marolf, Springer (2005),
\href{https://arxiv.org/abs/gr-qc/0309009}{arXiv:gr-qc/0309009}.

\bibitem{vanSuijlekom:2014pya}
W.~D.~van Suijlekom,
\emph{Noncommutative Geometry and Particle Physics},
Springer (2015),
\href{https://doi.org/10.1007/978-94-017-9162-5}{doi:10.1007/978-94-017-9162-5}.

\end{thebibliography}

\end{document}
